\documentclass[11pt]{article}
% ============================================================================
% _estilo.tex — sistema de diseño compartido para los PDF del frente
% Atención Armónica (Phideus). Derivado del sistema usado en el informe
% narrado de la memoria colectiva: paleta teal/dorado, cajas tcolorbox,
% tablas booktabs, más un vocabulario TikZ para los diagramas.
%
% Uso:  % ============================================================================
% _estilo.tex — sistema de diseño compartido para los PDF del frente
% Atención Armónica (Phideus). Derivado del sistema usado en el informe
% narrado de la memoria colectiva: paleta teal/dorado, cajas tcolorbox,
% tablas booktabs, más un vocabulario TikZ para los diagramas.
%
% Uso:  % ============================================================================
% _estilo.tex — sistema de diseño compartido para los PDF del frente
% Atención Armónica (Phideus). Derivado del sistema usado en el informe
% narrado de la memoria colectiva: paleta teal/dorado, cajas tcolorbox,
% tablas booktabs, más un vocabulario TikZ para los diagramas.
%
% Uso:  \input{_estilo}  justo antes de \begin{document}
%       \titleblock{Título}{Subtítulo}
% Compilar dos veces:  pdflatex -interaction=nonstopmode NOMBRE.tex
% ============================================================================

\usepackage[utf8]{inputenc}
\usepackage[T1]{fontenc}
\usepackage[spanish,es-noshorthands]{babel}
\usepackage{lmodern}
\usepackage[a4paper,margin=2.3cm]{geometry}
\usepackage{microtype}

\usepackage{xcolor}
\usepackage{amsmath}
\usepackage{amssymb}
\usepackage{booktabs}
\usepackage{colortbl}
\usepackage{array}
\usepackage{tabularx}
\usepackage{float}
\usepackage[most]{tcolorbox}
\usepackage{titlesec}
\usepackage{fancyhdr}
\usepackage{enumitem}
\usepackage{tikz}
\usetikzlibrary{positioning,arrows.meta,fit,backgrounds,calc,shapes.geometric,shapes.misc}
\usepackage[hidelinks]{hyperref}

% ---- Paleta ---------------------------------------------------------------
\definecolor{accent}{RGB}{20,105,95}      % teal profundo
\definecolor{accent2}{RGB}{196,142,53}    % dorado
\definecolor{inkgray}{RGB}{90,90,90}
\definecolor{rowgray}{RGB}{245,247,246}
\definecolor{alertred}{RGB}{160,58,48}

% ---- Secciones ------------------------------------------------------------
\titleformat{\section}
  {\sffamily\Large\bfseries\color{accent}}{\thesection}{0.6em}{}
  [{\color{accent2}\titlerule[1.4pt]}]
\titleformat{\subsection}
  {\sffamily\large\bfseries\color{accent}}{\thesubsection}{0.5em}{}
\titlespacing*{\section}{0pt}{1.4\baselineskip}{0.6\baselineskip}

% ---- Pie ------------------------------------------------------------------
\newcommand{\footername}{}
\pagestyle{fancy}
\fancyhf{}
\renewcommand{\headrulewidth}{0pt}
\renewcommand{\footrulewidth}{0.4pt}
\renewcommand{\footrule}{\hbox to\headwidth{\color{accent}\leaders\hrule height \footrulewidth\hfill}}
\fancyfoot[L]{\footnotesize\color{inkgray}\sffamily \footername}
\fancyfoot[R]{\footnotesize\color{inkgray}\sffamily \thepage}

% ---- Cajas ----------------------------------------------------------------
\newtcolorbox{keybox}[1]{
  enhanced, breakable, colback=accent!7, colframe=accent,
  boxrule=0.8pt, arc=2pt, left=8pt, right=8pt, top=6pt, bottom=6pt,
  coltitle=white, fonttitle=\sffamily\bfseries, title={#1}}

\newtcolorbox{alertbox}[1]{
  enhanced, breakable, colback=alertred!6, colframe=alertred,
  boxrule=0.8pt, arc=2pt, left=8pt, right=8pt, top=6pt, bottom=6pt,
  coltitle=white, fonttitle=\sffamily\bfseries, title={#1}}

\newtcolorbox{quietbox}{
  enhanced, breakable, colback=rowgray, colframe=rowgray,
  boxrule=0pt, arc=1pt, left=8pt, right=8pt, top=5pt, bottom=5pt,
  borderline west={2pt}{0pt}{accent!55}}

\newtcolorbox{goldbox}{
  enhanced, breakable, colback=accent2!8, colframe=accent2!8,
  boxrule=0pt, arc=1pt, left=8pt, right=8pt, top=5pt, bottom=5pt,
  borderline west={2pt}{0pt}{accent2}}

\newcommand{\headrow}{\rowcolor{accent!12}}
\setlist{leftmargin=1.4em, topsep=2pt, itemsep=1pt}
\renewcommand{\arraystretch}{1.15}

% ---- Bloque de título ------------------------------------------------------
\newcommand{\titleblock}[2]{%
  \noindent{\color{accent2}\rule{\linewidth}{1.6pt}}\par\smallskip
  {\sffamily\bfseries\Huge\color{accent} #1\par}\smallskip
  \ifx\relax#2\relax\else{\sffamily\large\color{inkgray} #2\par}\fi
  \smallskip\noindent{\color{accent2}\rule{\linewidth}{1.6pt}}\par\bigskip
  \gdef\footername{#1}%
}

% ---- Vocabulario de figuras TikZ ------------------------------------------
% Los diagramas ASCII de los originales se redibujan con estos estilos.
\tikzset{
  fbox/.style   ={draw=accent!65, fill=accent!8,  rounded corners=3pt, align=center,
                  inner sep=6pt, font=\sffamily\small, minimum height=8mm},
  fgold/.style  ={draw=accent2!85, fill=accent2!12, rounded corners=3pt, align=center,
                  inner sep=6pt, font=\sffamily\small, minimum height=8mm},
  fbad/.style   ={draw=alertred!70, fill=alertred!7, rounded corners=3pt, align=center,
                  inner sep=6pt, font=\sffamily\small, minimum height=8mm},
  fplain/.style ={draw=inkgray!45, fill=rowgray, rounded corners=3pt, align=center,
                  inner sep=6pt, font=\sffamily\small, minimum height=8mm},
  fdot/.style   ={circle, draw=accent!70, fill=accent!12, inner sep=0pt,
                  minimum size=6.5mm, font=\sffamily\scriptsize},
  fdot2/.style  ={circle, draw=accent2!90, fill=accent2!15, inner sep=0pt,
                  minimum size=6.5mm, font=\sffamily\scriptsize},
  flab/.style   ={font=\sffamily\scriptsize\color{inkgray}, align=center, inner sep=2pt},
  farr/.style   ={-{Stealth[length=2.6mm]}, draw=accent, thick},
  fgarr/.style  ={-{Stealth[length=2.6mm]}, draw=accent2!90!black, thick},
  fbarr/.style  ={-{Stealth[length=2.6mm]}, draw=alertred, thick},
  fdash/.style  ={-{Stealth[length=2.6mm]}, draw=inkgray!70, thick, dashed},
  fline/.style  ={draw=accent!55, thick},
  fxline/.style ={draw=alertred!80, thick, dash pattern=on 3pt off 2pt},
}

% Pie de figura, sin depender del paquete caption
\newcommand{\figcap}[1]{\par\smallskip{\sffamily\footnotesize\color{inkgray} #1\par}}

% Caja monoespaciada — SOLO para contenido ASCII puro (sin dibujo de cajas)
\newtcolorbox{codebox}{enhanced, breakable, colback=rowgray, colframe=accent!25,
  boxrule=0.5pt, arc=2pt, left=6pt, right=6pt, top=4pt, bottom=4pt,
  fontupper=\ttfamily\scriptsize}
  justo antes de \begin{document}
%       \titleblock{Título}{Subtítulo}
% Compilar dos veces:  pdflatex -interaction=nonstopmode NOMBRE.tex
% ============================================================================

\usepackage[utf8]{inputenc}
\usepackage[T1]{fontenc}
\usepackage[spanish,es-noshorthands]{babel}
\usepackage{lmodern}
\usepackage[a4paper,margin=2.3cm]{geometry}
\usepackage{microtype}

\usepackage{xcolor}
\usepackage{amsmath}
\usepackage{amssymb}
\usepackage{booktabs}
\usepackage{colortbl}
\usepackage{array}
\usepackage{tabularx}
\usepackage{float}
\usepackage[most]{tcolorbox}
\usepackage{titlesec}
\usepackage{fancyhdr}
\usepackage{enumitem}
\usepackage{tikz}
\usetikzlibrary{positioning,arrows.meta,fit,backgrounds,calc,shapes.geometric,shapes.misc}
\usepackage[hidelinks]{hyperref}

% ---- Paleta ---------------------------------------------------------------
\definecolor{accent}{RGB}{20,105,95}      % teal profundo
\definecolor{accent2}{RGB}{196,142,53}    % dorado
\definecolor{inkgray}{RGB}{90,90,90}
\definecolor{rowgray}{RGB}{245,247,246}
\definecolor{alertred}{RGB}{160,58,48}

% ---- Secciones ------------------------------------------------------------
\titleformat{\section}
  {\sffamily\Large\bfseries\color{accent}}{\thesection}{0.6em}{}
  [{\color{accent2}\titlerule[1.4pt]}]
\titleformat{\subsection}
  {\sffamily\large\bfseries\color{accent}}{\thesubsection}{0.5em}{}
\titlespacing*{\section}{0pt}{1.4\baselineskip}{0.6\baselineskip}

% ---- Pie ------------------------------------------------------------------
\newcommand{\footername}{}
\pagestyle{fancy}
\fancyhf{}
\renewcommand{\headrulewidth}{0pt}
\renewcommand{\footrulewidth}{0.4pt}
\renewcommand{\footrule}{\hbox to\headwidth{\color{accent}\leaders\hrule height \footrulewidth\hfill}}
\fancyfoot[L]{\footnotesize\color{inkgray}\sffamily \footername}
\fancyfoot[R]{\footnotesize\color{inkgray}\sffamily \thepage}

% ---- Cajas ----------------------------------------------------------------
\newtcolorbox{keybox}[1]{
  enhanced, breakable, colback=accent!7, colframe=accent,
  boxrule=0.8pt, arc=2pt, left=8pt, right=8pt, top=6pt, bottom=6pt,
  coltitle=white, fonttitle=\sffamily\bfseries, title={#1}}

\newtcolorbox{alertbox}[1]{
  enhanced, breakable, colback=alertred!6, colframe=alertred,
  boxrule=0.8pt, arc=2pt, left=8pt, right=8pt, top=6pt, bottom=6pt,
  coltitle=white, fonttitle=\sffamily\bfseries, title={#1}}

\newtcolorbox{quietbox}{
  enhanced, breakable, colback=rowgray, colframe=rowgray,
  boxrule=0pt, arc=1pt, left=8pt, right=8pt, top=5pt, bottom=5pt,
  borderline west={2pt}{0pt}{accent!55}}

\newtcolorbox{goldbox}{
  enhanced, breakable, colback=accent2!8, colframe=accent2!8,
  boxrule=0pt, arc=1pt, left=8pt, right=8pt, top=5pt, bottom=5pt,
  borderline west={2pt}{0pt}{accent2}}

\newcommand{\headrow}{\rowcolor{accent!12}}
\setlist{leftmargin=1.4em, topsep=2pt, itemsep=1pt}
\renewcommand{\arraystretch}{1.15}

% ---- Bloque de título ------------------------------------------------------
\newcommand{\titleblock}[2]{%
  \noindent{\color{accent2}\rule{\linewidth}{1.6pt}}\par\smallskip
  {\sffamily\bfseries\Huge\color{accent} #1\par}\smallskip
  \ifx\relax#2\relax\else{\sffamily\large\color{inkgray} #2\par}\fi
  \smallskip\noindent{\color{accent2}\rule{\linewidth}{1.6pt}}\par\bigskip
  \gdef\footername{#1}%
}

% ---- Vocabulario de figuras TikZ ------------------------------------------
% Los diagramas ASCII de los originales se redibujan con estos estilos.
\tikzset{
  fbox/.style   ={draw=accent!65, fill=accent!8,  rounded corners=3pt, align=center,
                  inner sep=6pt, font=\sffamily\small, minimum height=8mm},
  fgold/.style  ={draw=accent2!85, fill=accent2!12, rounded corners=3pt, align=center,
                  inner sep=6pt, font=\sffamily\small, minimum height=8mm},
  fbad/.style   ={draw=alertred!70, fill=alertred!7, rounded corners=3pt, align=center,
                  inner sep=6pt, font=\sffamily\small, minimum height=8mm},
  fplain/.style ={draw=inkgray!45, fill=rowgray, rounded corners=3pt, align=center,
                  inner sep=6pt, font=\sffamily\small, minimum height=8mm},
  fdot/.style   ={circle, draw=accent!70, fill=accent!12, inner sep=0pt,
                  minimum size=6.5mm, font=\sffamily\scriptsize},
  fdot2/.style  ={circle, draw=accent2!90, fill=accent2!15, inner sep=0pt,
                  minimum size=6.5mm, font=\sffamily\scriptsize},
  flab/.style   ={font=\sffamily\scriptsize\color{inkgray}, align=center, inner sep=2pt},
  farr/.style   ={-{Stealth[length=2.6mm]}, draw=accent, thick},
  fgarr/.style  ={-{Stealth[length=2.6mm]}, draw=accent2!90!black, thick},
  fbarr/.style  ={-{Stealth[length=2.6mm]}, draw=alertred, thick},
  fdash/.style  ={-{Stealth[length=2.6mm]}, draw=inkgray!70, thick, dashed},
  fline/.style  ={draw=accent!55, thick},
  fxline/.style ={draw=alertred!80, thick, dash pattern=on 3pt off 2pt},
}

% Pie de figura, sin depender del paquete caption
\newcommand{\figcap}[1]{\par\smallskip{\sffamily\footnotesize\color{inkgray} #1\par}}

% Caja monoespaciada — SOLO para contenido ASCII puro (sin dibujo de cajas)
\newtcolorbox{codebox}{enhanced, breakable, colback=rowgray, colframe=accent!25,
  boxrule=0.5pt, arc=2pt, left=6pt, right=6pt, top=4pt, bottom=4pt,
  fontupper=\ttfamily\scriptsize}
  justo antes de \begin{document}
%       \titleblock{Título}{Subtítulo}
% Compilar dos veces:  pdflatex -interaction=nonstopmode NOMBRE.tex
% ============================================================================

\usepackage[utf8]{inputenc}
\usepackage[T1]{fontenc}
\usepackage[spanish,es-noshorthands]{babel}
\usepackage{lmodern}
\usepackage[a4paper,margin=2.3cm]{geometry}
\usepackage{microtype}

\usepackage{xcolor}
\usepackage{amsmath}
\usepackage{amssymb}
\usepackage{booktabs}
\usepackage{colortbl}
\usepackage{array}
\usepackage{tabularx}
\usepackage{float}
\usepackage[most]{tcolorbox}
\usepackage{titlesec}
\usepackage{fancyhdr}
\usepackage{enumitem}
\usepackage{tikz}
\usetikzlibrary{positioning,arrows.meta,fit,backgrounds,calc,shapes.geometric,shapes.misc}
\usepackage[hidelinks]{hyperref}

% ---- Paleta ---------------------------------------------------------------
\definecolor{accent}{RGB}{20,105,95}      % teal profundo
\definecolor{accent2}{RGB}{196,142,53}    % dorado
\definecolor{inkgray}{RGB}{90,90,90}
\definecolor{rowgray}{RGB}{245,247,246}
\definecolor{alertred}{RGB}{160,58,48}

% ---- Secciones ------------------------------------------------------------
\titleformat{\section}
  {\sffamily\Large\bfseries\color{accent}}{\thesection}{0.6em}{}
  [{\color{accent2}\titlerule[1.4pt]}]
\titleformat{\subsection}
  {\sffamily\large\bfseries\color{accent}}{\thesubsection}{0.5em}{}
\titlespacing*{\section}{0pt}{1.4\baselineskip}{0.6\baselineskip}

% ---- Pie ------------------------------------------------------------------
\newcommand{\footername}{}
\pagestyle{fancy}
\fancyhf{}
\renewcommand{\headrulewidth}{0pt}
\renewcommand{\footrulewidth}{0.4pt}
\renewcommand{\footrule}{\hbox to\headwidth{\color{accent}\leaders\hrule height \footrulewidth\hfill}}
\fancyfoot[L]{\footnotesize\color{inkgray}\sffamily \footername}
\fancyfoot[R]{\footnotesize\color{inkgray}\sffamily \thepage}

% ---- Cajas ----------------------------------------------------------------
\newtcolorbox{keybox}[1]{
  enhanced, breakable, colback=accent!7, colframe=accent,
  boxrule=0.8pt, arc=2pt, left=8pt, right=8pt, top=6pt, bottom=6pt,
  coltitle=white, fonttitle=\sffamily\bfseries, title={#1}}

\newtcolorbox{alertbox}[1]{
  enhanced, breakable, colback=alertred!6, colframe=alertred,
  boxrule=0.8pt, arc=2pt, left=8pt, right=8pt, top=6pt, bottom=6pt,
  coltitle=white, fonttitle=\sffamily\bfseries, title={#1}}

\newtcolorbox{quietbox}{
  enhanced, breakable, colback=rowgray, colframe=rowgray,
  boxrule=0pt, arc=1pt, left=8pt, right=8pt, top=5pt, bottom=5pt,
  borderline west={2pt}{0pt}{accent!55}}

\newtcolorbox{goldbox}{
  enhanced, breakable, colback=accent2!8, colframe=accent2!8,
  boxrule=0pt, arc=1pt, left=8pt, right=8pt, top=5pt, bottom=5pt,
  borderline west={2pt}{0pt}{accent2}}

\newcommand{\headrow}{\rowcolor{accent!12}}
\setlist{leftmargin=1.4em, topsep=2pt, itemsep=1pt}
\renewcommand{\arraystretch}{1.15}

% ---- Bloque de título ------------------------------------------------------
\newcommand{\titleblock}[2]{%
  \noindent{\color{accent2}\rule{\linewidth}{1.6pt}}\par\smallskip
  {\sffamily\bfseries\Huge\color{accent} #1\par}\smallskip
  \ifx\relax#2\relax\else{\sffamily\large\color{inkgray} #2\par}\fi
  \smallskip\noindent{\color{accent2}\rule{\linewidth}{1.6pt}}\par\bigskip
  \gdef\footername{#1}%
}

% ---- Vocabulario de figuras TikZ ------------------------------------------
% Los diagramas ASCII de los originales se redibujan con estos estilos.
\tikzset{
  fbox/.style   ={draw=accent!65, fill=accent!8,  rounded corners=3pt, align=center,
                  inner sep=6pt, font=\sffamily\small, minimum height=8mm},
  fgold/.style  ={draw=accent2!85, fill=accent2!12, rounded corners=3pt, align=center,
                  inner sep=6pt, font=\sffamily\small, minimum height=8mm},
  fbad/.style   ={draw=alertred!70, fill=alertred!7, rounded corners=3pt, align=center,
                  inner sep=6pt, font=\sffamily\small, minimum height=8mm},
  fplain/.style ={draw=inkgray!45, fill=rowgray, rounded corners=3pt, align=center,
                  inner sep=6pt, font=\sffamily\small, minimum height=8mm},
  fdot/.style   ={circle, draw=accent!70, fill=accent!12, inner sep=0pt,
                  minimum size=6.5mm, font=\sffamily\scriptsize},
  fdot2/.style  ={circle, draw=accent2!90, fill=accent2!15, inner sep=0pt,
                  minimum size=6.5mm, font=\sffamily\scriptsize},
  flab/.style   ={font=\sffamily\scriptsize\color{inkgray}, align=center, inner sep=2pt},
  farr/.style   ={-{Stealth[length=2.6mm]}, draw=accent, thick},
  fgarr/.style  ={-{Stealth[length=2.6mm]}, draw=accent2!90!black, thick},
  fbarr/.style  ={-{Stealth[length=2.6mm]}, draw=alertred, thick},
  fdash/.style  ={-{Stealth[length=2.6mm]}, draw=inkgray!70, thick, dashed},
  fline/.style  ={draw=accent!55, thick},
  fxline/.style ={draw=alertred!80, thick, dash pattern=on 3pt off 2pt},
}

% Pie de figura, sin depender del paquete caption
\newcommand{\figcap}[1]{\par\smallskip{\sffamily\footnotesize\color{inkgray} #1\par}}

% Caja monoespaciada — SOLO para contenido ASCII puro (sin dibujo de cajas)
\newtcolorbox{codebox}{enhanced, breakable, colback=rowgray, colframe=accent!25,
  boxrule=0.5pt, arc=2pt, left=6pt, right=6pt, top=4pt, bottom=4pt,
  fontupper=\ttfamily\scriptsize}


\begin{document}

\titleblock{Atención armónica}{Aprender la estructura en lugar de medirla}

\begin{quietbox}
Documento de encuadre del frente de \emph{atención armónica} de Phideus: qué pregunta abre, qué tomó de AlphaFold y dónde esa analogía se rompe, cómo se vuelve medible el problema de repartir parciales por fuente, qué arquitecturas se comparan y qué decidió la \texttt{Fase 0}.
\end{quietbox}

\section{La pregunta que abre el frente}

Phideus viene sosteniendo una hipótesis precisa: las relaciones de frecuencia —los ratios entre los tonos que componen un sonido— forman un lenguaje informacional, y una red neuronal puede aprovecharlo. Hasta ahora ese lenguaje entró a los modelos como un descriptor calculado a mano e \emph{inyectado} en un encoder genérico: el modelo recibía, además del audio, una lista de números que resumían la estructura armónica, y el experimento medía si esa lista ayudaba. El frente de atención armónica cambia el lugar donde vive el conocimiento. En vez de calcular los ratios afuera y pasárselos al modelo, busca una arquitectura cuya atención opere directamente sobre la geometría armónica del sonido, de modo que la red infiera la estructura en lugar de recibirla resumida.

El cambio no es de grado sino de naturaleza. Un descriptor inyectado es una respuesta entregada; una geometría interna es un espacio donde la respuesta se construye. La pregunta del frente, entonces, es si una red puede razonar con ratios —usar la consistencia global de las relaciones de frecuencia para resolver lo que cada relación aislada deja ambiguo— o si el aporte de los ratios se agota en pasarlos como features.

\section{La lección de AlphaFold, bien entendida}

El disparador de este frente fue AlphaFold, el sistema que predice la forma tridimensional de una proteína a partir de su secuencia. Su potencia suele atribuirse a que usa transformers. Esa lectura se queda corta y conviene desarmarla, porque lo que importa para Phideus es otra cosa.

Una proteína es una cadena de aminoácidos que se pliega en el espacio. Predecir el plegamiento es predecir, para cada par de aminoácidos, qué tan cerca quedan en la estructura final. AlphaFold no trata a los aminoácidos como una secuencia plana a la que aplica atención: construye una \textbf{representación de pares}, una tabla donde cada casilla $(i,j)$ acumula lo que el modelo cree sobre la relación entre el aminoácido $i$ y el $j$. Esa tabla es el objeto central del cómputo, y la operación que la hace funcionar es la \textbf{actualización triangular}: para revisar lo que sabe del par $(i,j)$, el modelo mira todos los terceros aminoácidos $k$ y combina lo que sabe de $(i,k)$ con lo que sabe de $(k,j)$. La intuición es geométrica y estricta: si A está cerca de B y B está cerca de C, eso restringe cuán lejos puede estar A de C. La desigualdad triangular del espacio físico se vuelve una operación aprendible sobre tripletes.

\begin{keybox}{La enseñanza transferible}
La inteligencia de AlphaFold no es la atención en abstracto, sino que la atención vive \textbf{dentro de la geometría natural del problema}: razona en el espacio de relaciones donde las reglas del plegamiento tienen forma, y deja que la consistencia global entre pares —la restricción triangular— corrija las estimaciones locales ruidosas. Una distancia entre dos residuos, tomada aislada, es una conjetura; el conjunto de todas las distancias tiene que ser realizable como una única figura en tres dimensiones, y esa exigencia de coherencia es la que convierte muchas conjeturas dudosas en una estructura.
\end{keybox}

\section{Dónde la analogía se rompe, y dónde se sostiene}

Trasladar esto a la música tiene una trampa que es necesario nombrar, porque cuesta cara si se pasa por alto.

\begin{alertbox}{La pieza que, traída literalmente, queda hueca}
La tentación inmediata es armar una actualización triangular sobre las frecuencias en escala logarítmica: si la distancia de A a B es $\log(f_B) - \log(f_A)$, y la de B a C es $\log(f_C) - \log(f_B)$, entonces la de A a C debería ser la suma. Pero esa ``restricción'' es una identidad algebraica: vale para tres números cualesquiera, siempre, exactamente. No hay nada que enforcar ni nada que aprender. En el espacio tridimensional la desigualdad triangular recorta el conjunto de configuraciones posibles; en la recta de las log-frecuencias la suma de diferencias no recorta nada.
\end{alertbox}

La salida es reubicar la no-trivialidad donde sí habita en el dominio armónico. Un sonido con altura definida está hecho de \textbf{parciales}: la frecuencia fundamental y una serie de componentes en frecuencias aproximadamente múltiplas de ella. Cuando suenan varias notas a la vez —una mezcla polifónica—, el oído recibe un solo conjunto de parciales entreverados y tiene que repartirlos: decidir qué parcial pertenece a qué fuente. Ese reparto es un problema de inferencia global. La relación que importa no es la distancia continua entre dos parciales, sino la pertenencia: \emph{¿estos dos parciales nacen de la misma fundamental?} Y la pertenencia a una misma fuente es una relación de equivalencia, lo que la vuelve no trivial: si el parcial $i$ pertenece a la misma fuente que $j$, y $j$ a la misma que $k$, entonces $i$ y $k$ tienen que pertenecer a la misma fuente. Esa transitividad es la análoga genuina de la restricción triangular. Aparece donde la evidencia de a pares es ambigua y solo la coherencia del conjunto resuelve.

El caso que vuelve tangible la ambigüedad es una mezcla de dos notas cuyas fundamentales están en relación simple. Tomemos fundamentales en 100 y 150 hertz.

\begin{figure}[H]
\centering
\begin{tikzpicture}[x=1.48cm, y=1cm]
  % lanes
  \node[flab, anchor=east] at (0.35, 2.35) {fuente A\\$f_0 = 100$};
  \node[flab, anchor=east] at (0.35, 1.15) {fuente B\\$f_0 = 150$};

  % shared verticals
  \draw[inkgray!40, dashed] (3.0, -0.55) -- (3.0, 2.85);
  \draw[inkgray!40, dashed] (6.0, -0.55) -- (6.0, 2.85);

  % axis
  \draw[farr] (0.35, 0) -- (6.9, 0);
  \node[flab, anchor=west] at (6.95, 0) {frecuencia (Hz)};

  % ticks
  \foreach \x/\lb in {1/100, 1.5/150, 2/200, 3/300, 4/400, 4.5/450, 5/500, 6/600} {
    \draw[inkgray!60] (\x, 0) -- (\x, -0.13);
    \node[flab, anchor=north, inner sep=1pt] at (\x, -0.16) {\lb};
  }

  % source A partials
  \foreach \x in {1, 2, 4, 5} { \node[fdot] at (\x, 2.35) {}; }
  \foreach \x in {3, 6} { \node[fdot2] at (\x, 2.35) {}; }
  \draw[fline] (1, 2.35) -- (6, 2.35);
  \foreach \x in {1, 2, 4, 5} { \node[fdot] at (\x, 2.35) {}; }
  \foreach \x in {3, 6} { \node[fdot2] at (\x, 2.35) {}; }
  \node[flab, anchor=west] at (6.32, 2.35) {\dots};

  % source B partials
  \draw[fline] (1.5, 1.15) -- (6, 1.15);
  \foreach \x in {1.5, 4.5} { \node[fdot] at (\x, 1.15) {}; }
  \foreach \x in {3, 6} { \node[fdot2] at (\x, 1.15) {}; }
  \node[flab, anchor=west] at (6.32, 1.15) {\dots};

  % annotations
  \node[flab, anchor=north] at (3.0, -0.62) {$300 = 3\times 100 = 2\times 150$};
  \node[flab, anchor=north] at (6.0, -0.62) {$600 = 6\times 100 = 4\times 150$};
\end{tikzpicture}
\figcap{Dos series de parciales sobre el eje de frecuencia. Los parciales en dorado (300 y 600 Hz) pertenecen coherentemente a cualquiera de las dos series: mirados aislados, no se pueden asignar.}
\end{figure}

El parcial en 300 hertz es, a la vez, el tercer armónico de A y el segundo de B. Mirado solo, no se puede asignar: pertenece coherentemente a cualquiera de las dos series. Lo que desambigua no es ese parcial sino el resto de la evidencia —que 200 y 400 anclan a A, que 450 ancla a B— y la exigencia de que la asignación final sea una partición consistente. Resolver esto pide integrar la estructura completa, no leer un par por vez. Ese es exactamente el régimen donde una representación de pares con propagación transitiva puede aportar sobre un modelo que clasifica cada par por separado.

\section{El problema concreto: repartir parciales por fuente}

El experimento decisivo del frente toma esa intuición y la vuelve medible en su forma más limpia. La tarea es el agrupamiento armónico: dado el conjunto de parciales de una mezcla polifónica, predecir la matriz de ``misma fuente'' —para cada par de parciales, si nacen o no de la misma fundamental. Esa matriz es el objeto central, el equivalente del mapa de contactos de AlphaFold.

La materia prima es audio sintético, generado de modo que se conozca con exactitud qué parcial vino de qué fuente. Esa exactitud es la condición que hace posible el experimento: AlphaFold pudo aprender porque tuvo decenas de miles de estructuras resueltas como verdad de referencia, y acá la verdad de referencia se fabrica al sintetizar la señal. Cada mezcla combina una, dos o tres fuentes; cada fuente aporta sus parciales con una fundamental propia; el agrupamiento verdadero queda registrado pieza por pieza. La supervisión es perfecta porque la construimos.

\section{Las arquitecturas que se comparan}

La pregunta del frente solo tiene respuesta si se aísla con cuidado qué se está midiendo. Por eso el experimento no enfrenta un modelo contra otro cualquiera, sino una escalera de arquitecturas donde cada peldaño agrega una sola cosa.

El primer peldaño es un transformer que mira los parciales como un conjunto de elementos y atiende entre ellos, con un sesgo de atención derivado de la distancia en log-frecuencia. Lee cada par desde las representaciones de los dos parciales. No recibe ninguna pista armónica explícita: tiene que descubrir la estructura desde la frecuencia y la amplitud crudas.

\begin{figure}[H]
\centering
\begin{tikzpicture}[node distance=8mm and 9mm]
  \node[fbox] (p) {parciales\\$[\log f,\ \log a]$};
  \node[fbox, right=of p] (emb) {embedding};
  \node[fbox, right=of emb] (att) {self-attention $\times L$\\{\scriptsize sesgo $=\Delta\log f$}};
  \node[fbox, below=12mm of att] (rd) {readout$(\text{token}_i,\ \text{token}_j)$};
  \node[fgold, left=12mm of rd] (out) {¿misma fuente?};
  \draw[farr] (p) -- (emb);
  \draw[farr] (emb) -- (att);
  \draw[farr] (att) -- (rd);
  \draw[farr] (rd) -- (out);
\end{tikzpicture}
\figcap{\texttt{A-naive}. Atención entre tokens de parcial; el par se lee desde las representaciones de los dos parciales. Sin pistas armónicas.}
\end{figure}

El segundo peldaño recibe, además, las relaciones armónicas calculadas a mano para cada par: qué tan cerca está su cociente de un ratio simple, qué fundamental común implícita comparten. Es el mismo transformer, pero con la pista entregada. Funciona como el mejor competidor posible \emph{a nivel de par}: si la respuesta ya está contenida en esas pistas, este modelo la alcanza.

\begin{figure}[H]
\centering
\begin{tikzpicture}[node distance=8mm and 9mm]
  \node[fbox] (p) {parciales};
  \node[fbox, right=of p] (emb) {embedding};
  \node[fbox, right=of emb] (att) {self-attention $\times L$};
  \node[fbox, below=14mm of att] (rd) {readout$(\text{token}_i,\ \text{token}_j,\ \text{pair-features})$};
  \node[fgold, left=10mm of rd, text width=3.1cm] (pf) {pair-features$(i,j)$\\{\scriptsize $\Delta\log f$, cercanía a ratio simple, fundamental común implícita, \dots}};
  \node[fgold, below=9mm of rd] (out) {¿misma fuente?};
  \draw[farr] (p) -- (emb);
  \draw[farr] (emb) -- (att);
  \draw[farr] (att) -- (rd);
  \draw[fgarr] (pf) -- (rd);
  \draw[farr] (rd) -- (out);
\end{tikzpicture}
\figcap{\texttt{A-rich}. El mismo transformer, con las relaciones armónicas del par entregadas al readout. Es el mejor competidor posible a nivel de par.}
\end{figure}

El tercer peldaño es la arquitectura que el frente pone a prueba. Mantiene una representación de pares que se actualiza a lo largo de las capas, deja que la atención entre parciales se sesgue por esa representación, y aplica la actualización triangular que propaga la transitividad: para revisar el par $(i,j)$, combina lo que sabe de $(i,k)$ y $(k,j)$ sobre todos los terceros $k$. La lectura final sale de la representación de pares, no de los tokens sueltos.

\begin{figure}[H]
\centering
\begin{tikzpicture}[node distance=7mm and 10mm]
  \node[fbox] (tok) {parciales $\to$ tokens};
  \node[fgold, right=14mm of tok, text width=3.5cm] (z) {pair-features $\to$ $z[i,j]$\\{\scriptsize representación de pares}};

  \node[fbox, below=13mm of tok, text width=3.4cm] (a1) {atención entre tokens\\sesgada por $z$};
  \node[fbox, below=of a1, text width=4.6cm] (a2) {$z \leftarrow z + \text{comunicación}(\text{token}_i,\text{token}_j)$};
  \node[fgold, below=of a2, text width=5.6cm] (a3) {\textsc{triángulo}\\[1pt]$z \leftarrow z + \sum_k a(z[i,k]) \odot b(z[k,j])$};

  \begin{scope}[on background layer]
    \node[draw=inkgray!45, dashed, rounded corners=4pt, fill=rowgray!60,
          inner sep=8pt, fit=(a1)(a2)(a3)] (loop) {};
  \end{scope}
  \node[flab, anchor=north west, font=\sffamily\scriptsize\bfseries\color{inkgray}]
        at ([xshift=2pt, yshift=-2pt] loop.north west) {repetir $L$ veces};

  \node[flab, right=3mm of a3, text width=2.4cm, anchor=west]
        {propaga\\transitividad};

  \node[fbox, below=11mm of loop] (rd) {readout$(z[i,j])$};
  \node[fgold, below=8mm of rd] (out) {¿misma fuente?};

  \draw[farr] (tok) -- (a1);
  \draw[fgarr] (z) |- (a1.east);
  \draw[farr] (a1) -- (a2);
  \draw[farr] (a2) -- (a3);
  \draw[farr] (loop.south) -- (rd);
  \draw[farr] (rd) -- (out);
\end{tikzpicture}
\figcap{\texttt{B} (\emph{Harmonic Pairformer}). La representación de pares $z$ persiste a lo largo de las capas, sesga la atención entre tokens y se actualiza con la suma sobre terceros. La lectura final sale de $z[i,j]$, no de los tokens sueltos.}
\end{figure}

La comparación que importa es \texttt{B} contra el segundo peldaño, con las pistas armónicas igualadas: si \texttt{B} gana, gana la \emph{maquinaria} —la representación de pares y la propagación transitiva— y no el hecho de tener mejores features. Para separar el aporte del triángulo del aporte de simplemente mantener una representación de pares, el experimento incluye una variante de \texttt{B} con la misma cantidad de parámetros pero sin la suma sobre terceros: si \texttt{B} le gana también a esa variante, lo que aporta es específicamente la transitividad. Un control adicional, con las pistas armónicas barajadas, vigila que cualquier ventaja no venga del mero tamaño del modelo.

\section{La síntesis no alcanza con generarla: hay que probar que deja lugar}

El punto más delicado del frente no es la arquitectura sino la validez del experimento, y conviene exponerlo porque enseña algo que excede a este caso. Un experimento que compara \texttt{B} contra el competidor con features igualadas solo dice algo si las features no resuelven ya la tarea por sí solas. Si las pistas armónicas entregadas alcanzan para separar las fuentes, el competidor llega al techo y \texttt{B} no tiene dónde mejorar: una eventual igualdad entre ambos no probaría que la transitividad no sirve, sino que la tarea no dejaba lugar para que sirviera.

Esa trampa apareció dos veces, y las dos veces una auditoría previa la atrapó antes de gastar cómputo. En el primer diseño, con parciales en múltiplos enteros exactos, dos parciales de la misma fuente quedaban en un cociente de enteros pequeños —un ratio simple— y una sola feature calculada los separaba casi perfectamente: la tarea era trivial a nivel de par. En el segundo diseño, con los parciales ya desafinados para romper ese atajo, sobrevivía otro: la envolvente de amplitud, que decaía de forma fija con el número de armónico, filtraba indirectamente cuál era ese número, y volvía a hacer separable cada par. Ninguno de los dos atajos era visible a simple vista; ambos habrían producido un resultado falsamente nulo, leído como ``la transitividad no aporta'' cuando en realidad el dato no la necesitaba.

\begin{keybox}{La norma que queda}
Ningún conjunto de datos pasa a entrenamiento sin demostrar dos cosas a la vez. Que ninguna combinación cerrada de las pistas disponibles resuelva la tarea por sí sola —hay margen genuino para que la maquinaria global aporte—, y que la tarea siga siendo resoluble cuando se la mira en conjunto —no es imposible, solo localmente ambigua. El experimento decisivo vale en la medida en que existe esa franja entre lo trivial y lo imposible, y la franja se verifica antes, no después.
\end{keybox}

\section{Qué decidió Fase 0}

El frente no prometía que una red razonara con armonía; puso esa afirmación en riesgo de ser refutada barata. La \texttt{Fase 0} ya produjo una primera respuesta matizada. El salto más fuerte no vino de una regla local de ratios ni de una feature externa, sino de sostener una representación explícita de pares: \texttt{B-minus} $\gg$ \texttt{A-rich} mostró que el objeto útil del cómputo no es solo el pico, sino la relación entre picos.

El \texttt{triangle}, en cambio, no quedó validado como ganador universal. En \texttt{IID} y \texttt{OOD-regime}, una mezcla local param-matched puede igualarlo o superarlo levemente. Pero en \texttt{OOD-poly}, donde la polifonía del test aumenta y la ambigüedad global se vuelve más dura, \texttt{B} supera a \texttt{B-local} en \texttt{AUC/AP} y también supera claramente a \texttt{B-shuffle}. Esa combinación sostiene una lectura acotada: la estructura triangular no es simple capacidad adicional, sino un sesgo relacional que ayuda a generalizar cuando cambia la cantidad de fuentes.

\begin{keybox}{El caveat es parte del resultado}
La red puede rankear mejor los pares y, sin embargo, fallar como sistema de agrupamiento si el umbral $\tau$ elegido en validación no transfiere. Por eso la salida de \texttt{Fase 0} no es una promoción sin reservas, sino un \texttt{GO} acotado: la arquitectura tiene señal, pero hay que resolver la decisión de clustering.
\end{keybox}

\section{Geometría relacional de la armonía}

La intuición geométrica del frente no debe confundirse con una geometría euclídea de frecuencias. En $\log f$, las diferencias entre tres picos cumplen identidades algebraicas triviales; esa no es una restricción aprendible. La geometría que sí aparece es relacional.

\begin{table}[H]
\centering
\begin{tabular}{ll}
\toprule
\headrow \textbf{Objeto del dominio} & \textbf{Objeto de la geometría relacional} \\
\midrule
picos espectrales & nodos \\
\texttt{same-source[i,j]} & aristas aprendidas \\
fuentes armónicas & clases de equivalencia \\
\bottomrule
\end{tabular}
\figcap{El mapeo entre el dominio armónico y el grafo sobre el que opera la red.}
\end{table}

Cada fuente puede pensarse como una familia generativa discreta, por ejemplo:
\[
  f_n = n\,f_0\sqrt{1 + \beta n^2}
\]

Los picos de una misma fuente no son solo cercanos; son coherentes con un mismo generador latente. La matriz de pares \texttt{z[i,j]} es el espacio donde la red representa esa compatibilidad. La actualización triangular propaga evidencia indirecta: si $i$ parece ir con $k$, y $k$ con $j$, entonces $i$ y $j$ reciben información nueva. No se impone transitividad como regla lógica dura; se aprende cuándo esa evidencia global corrige la ambigüedad local.

Esa es la diferencia con un descriptor clásico. Un descriptor calcula afuera una señal armónica y se la entrega al modelo. El \emph{Harmonic Pairformer} intenta que la relación armónica sea el medio interno del cómputo.

\section{Qué sigue}

El siguiente paso ya no es una \texttt{Fase 0.5} abstracta de ``calibración de $\tau$'', porque esa auditoría ya se hizo y dejó una corrección metodológica más fuerte: el problema de \texttt{B} en \texttt{OOD-poly} no estaba en $\tau$, sino en \texttt{connected-components} como lector de la matriz de pares. La ventaja representacional de \texttt{B} sí aparece; y \texttt{Fase 0.6} ya agregó la pieza que faltaba para decir algo más fuerte: clusterers globales deployables (\texttt{spectral} y \texttt{agglo} con \texttt{k} estimado) ya recuperan una ventaja real de \texttt{B} sobre \texttt{B-local} en \texttt{OOD-poly}. Lo que falta ahora no es un umbral mejor, sino cerrar el gap que deja la subestimación de \texttt{k}.

Por eso, antes de saltar a audio real, el paso más razonable ya no es ``probar algún clusterer global'' en abstracto, porque eso ya ocurrió. La decisión real es si conviene abrir un \textbf{Stage B} donde cambie no la red pairwise sino la lectura de su salida mediante una cabeza explícita de \texttt{k}/partición, o si se pasa directo a \texttt{Fase 1a}: renderizar las mezclas sintéticas, detectar picos con CQT y usar esos picos detectados como tokens. Ahí entran picos faltantes, picos corridos, picos fusionados y espurios, pero todavía con ground truth controlado. Audio real/stems queda como una fase posterior.

\end{document}
